\documentclass[a4paper]{article}
\usepackage[pdftex]{graphicx}
\usepackage[utf8]{inputenc}
\usepackage{enumerate}
\usepackage{icomma}
\usepackage{siunitx}
\sisetup{locale=DE} 
\usepackage{tikz}
\usepackage{href-ul}
\hypersetup{
	colorlinks=true,
	linkcolor=blue,
	urlcolor=blue}
\usepackage{amssymb}
\usepackage{geometry}
\geometry{a4paper, top=15mm, left=15mm, right=15mm, bottom=15mm,
	headsep=10mm, footskip=12mm}

\begin{document}
\begin{tikzpicture}(10,3)
	\draw[ultra thick](2,0) --(10,0) -- (10,1.3) --(2,1.3) -- (2,0);
	\draw[fill=black](2,0)-- (0,.6) -- (2,.6) -- (2,0);
	\node at (6,.6) {\large Lösungsblatt von \href{https://www.okuyakl.de}{www.okuyakl.de}};
\end{tikzpicture}
\vspace{0.5 cm}

\noindent
\begin{minipage}{0.5\textwidth}
\noindent{\bf Aufgabe 1.}\\
Der Funktionsterm ist 
$$y =5^x$$

\noindent{\bf Aufgabe 2.}\\
$$ y = 100 \cdot 0,8^x$$


\end{minipage}
\hspace{0.5 cm}
\begin{minipage}{0.5\textwidth}
\noindent{\bf Aufgabe 3. a)}\\
\includegraphics[width=7 cm]{exdrei029}
\end{minipage}

\noindent{\bf Aufgabe 3. a)}\\
$$
\renewcommand{\arraystretch}{2}
\begin{array}{c|cccccccccc}
x & 1 & 2 & 3 & 4 & 5 & 6 & 7 & 8 & 9 & 10 \\
\hline
y & 1,58 & 2,5 & 3,95 & 6,25 & 9,88 & 15,63 & 24,71 & 39,06 & 61,76 & 97,66 \\
\end{array}
$$	

\noindent{\bf Aufgabe 3 b.)}\\
Den Wachstumsfaktor berechnen wir mit 
$$a = 2,5^{0,5} = 1,58 $$
Die Zunahme beträgt damit 58\%.

\noindent{\bf Aufgabe 3. c)}\\ 
Wir lesen ab: 50 Erkrankte gibt es nach 8,5 Tagen.

\fbox{
\noindent
\begin{minipage}{0.5\textwidth}
	\noindent{\bf Aufgabe 4. b)}\\
	Den Wachstumsfaktor berechnen wir mit 
	$$a = 3^{0,5} = 1,73 $$
	Die Zunahme beträgt damit 73\%.
	
	\noindent{\bf Aufgabe 4. c)}\\
	Wir lesen ab: 100 Bakterien gibt es nach 4,2 Stunden.
\end{minipage}
\hfill
\begin{minipage}{0.4\textwidth}
	\noindent{\bf Aufgabe 4. a)}\\
	\includegraphics[width=7 cm]{exvier029}
\end{minipage}}

\noindent{\bf Aufgabe 5. a)}\\
Der Funktionsterm lautet:
$$y = 0,5 \cdot 1,35^x$$

\noindent{\bf Aufgabe 5. b)}\\
$$
\renewcommand{\arraystretch}{2}
\begin{array}{c|ccccccccc}
x & 0 & 1 & 2 & 3 & 4 & 5 & 6 & 7 & 8 \\
\hline
y & 0,50 & 0,67 & 0,91 & 1,23 & 1,66 & 2,24 & 3,03 & 4,09 & 5,51
\end{array}
$$
\noindent
\begin{minipage}{0.5\textwidth}
	\noindent{\bf Aufgabe 5. b)}\\
	\includegraphics[width=7 cm]{exfunf029}
\end{minipage}
\hfill
\begin{minipage}{0.5\textwidth}
   \noindent{\bf Aufgabe 5. c)}\\
   Ein Fünftel der Wasserfläche sind $\SI{4}{\meter^2}$. Diese werden nach 7 Tagen erreicht, also am 17. Juni.
\end{minipage}

\noindent
\begin{minipage}{0.5\textwidth}
	\noindent{\bf Aufgabe 6. a)}\\
	\includegraphics[width=7 cm]{exse029}
\end{minipage}
\hfill
\begin{minipage}{0.5\textwidth}
		\noindent{\bf Aufgabe 6. b)}\\
		Den Wachstumsfaktor berechnen wir mit 
		$$a = 0,5^{0,5} = 0,71 $$
		Die Abnahme beträgt damit $100\%-71\%=29\%$.
		
		\noindent{\bf Aufgabe 6. c)}\\
		Wir lesen ab: nach 1,5 Tagen sind noch $\SI{3}{\milli\gram}$ vorhanden.
\end{minipage}
\vspace{0.5 cm}

\noindent{\bf Aufgabe 7.}\\
Von 2016 bis 2028 sind es 12 Jahre. Der Wachstumsfaktor ist 
$$a=(1-{p \over 100})= 1-0,0125 = 0,9875$$
Der Schadstoffausstoß nach 12 Jahren ist:
$$y=100\% \cdot 0,9875^{12} =85,99\% $$
Er hat also um etwas über 14 \% abgenommen. Das Ziel wird erreicht.

\begin{center}
	\includegraphics[width=7 cm]{../../viecher/endcomic.pdf}
	
	Hier geht es zurück zum \href{https://www.okuyakl.de/math/m10expL029/aa029.pdf}{Aufgabenblatt}
\end{center}

\end{document}

